\begin{tomtat}
	\textbf{ĐỀ 03}
\end{tomtat}
\begin{ex}
	Giải hệ phương trình $\begin{cases} 3x - 2y = 9 \\ x + 3y = -8 \end{cases}$
	\loigiai{
		Ta có hệ phương trình:
		$$ \begin{cases} 3x - 2y = 9 & (1) \\ x + 3y = -8 & (2) \end{cases} $$
		Từ phương trình (2), ta có thể biểu diễn $x$ theo $y$:
		$x = -8 - 3y$.
		Thay biểu thức của $x$ này vào phương trình (1):
		\begin{align*}
			3(-8 - 3y) - 2y &= 9 \\
			-24 - 9y - 2y &= 9 \\
			-11y &= 9 + 24 \\
			-11y &= 33 \\
			y &= \frac{33}{-11} \\
			y &= -3
		\end{align*}
		Thay giá trị $y=-3$ trở lại vào biểu thức $x = -8 - 3y$:
		\begin{align*}
			x &= -8 - 3(-3) \\
			x &= -8 + 9 \\
			x &= 1
		\end{align*}
		Vậy nghiệm của hệ phương trình là $(x; y) = (1; -3)$.
	}
\end{ex}


\begin{ex}
	Rút gọn biểu thức $B = \left(2 + \frac{a+3\sqrt{a}}{\sqrt{a}+3}\right)\left(2 + \frac{a-3\sqrt{a}}{3-\sqrt{a}}\right)$, với $a \ge 0, a \ne 9$.
	\loigiai{
		Với điều kiện $a \ge 0, a \ne 9$, ta có:
		\begin{align*}
			B &= \left(2 + \frac{\sqrt{a}(\sqrt{a}+3)}{\sqrt{a}+3}\right)\left(2 + \frac{\sqrt{a}(\sqrt{a}-3)}{-(\sqrt{a}-3)}\right) \\
			&= \left(2 + \sqrt{a}\right)\left(2 - \sqrt{a}\right) \\
			&= 2^2 - (\sqrt{a})^2 \\
			&= 4 - a
		\end{align*}
		Vậy biểu thức rút gọn là $B = 4-a$.
	}
\end{ex}


\begin{ex}
	Gọi $x_1, x_2$ là hai nghiệm của phương trình $x^2 - 7x + 4 = 0$. Không giải phương trình hãy tính giá trị của biểu thức $P = x_1 + x_2 - 3x_1x_2$.
	\loigiai{
		Phương trình đã cho là $x^2 - 7x + 4 = 0$.
		Đây là phương trình bậc hai có dạng $ax^2+bx+c=0$ với $a=1$, $b=-7$, $c=4$.
		Ta tính biệt thức $\Delta = b^2-4ac = (-7)^2 - 4(1)(4) = 49 - 16 = 33$.
		Vì $\Delta = 33 > 0$ nên phương trình có hai nghiệm phân biệt $x_1, x_2$.
		
		Theo hệ thức Vi-ét, ta có:
		$$ \begin{cases} 
			x_1+x_2 = \dfrac{-b}{a} = \dfrac{-(-7)}{1} = 7 \\ 
			x_1x_2 = \dfrac{c}{a} = \dfrac{4}{1} = 4 
		\end{cases} $$
		Ta cần tính giá trị biểu thức $P = x_1 + x_2 - 3x_1x_2$.
		Thay giá trị $x_1+x_2$ và $x_1x_2$ đã tính ở trên vào biểu thức $P$:
		\begin{align*}
			P &= (x_1+x_2) - 3(x_1x_2) \\
			&= (7) - 3(4) \\
			&= 7 - 12 \\
			&= -5
		\end{align*}
		Vậy giá trị của biểu thức $P$ là -5.
	}
\end{ex}


\begin{ex}
	Một cái hộp đựng 3 viên bi đỏ và 2 viên bi xanh. Lấy lần lượt 2 viên bi từ cái hộp đó (không hoàn lại). Tính xác suất để lấy được 2 viên bi là bi xanh.
	\loigiai{
		Tổng số viên bi trong hộp là: $3 \text{ (đỏ)} + 2 \text{ (xanh)} = 5$ viên bi.
		Ký hiệu 3 viên bi đỏ là Đ$_1$, Đ$_2$, Đ$_3$ và 2 viên bi xanh là X$_1$, X$_2$.
		Ta lấy lần lượt 2 viên bi không hoàn lại. Mỗi kết quả là một cặp sắp thứ tự (viên thứ nhất, viên thứ hai).
		Ta cần tính xác suất để cả hai viên bi lấy được đều là bi xanh.
		
		\textbf{Phương pháp: Liệt kê và đếm số trường hợp}
		\begin{itemize}
			\item \textbf{Không gian mẫu $\Omega$ (tất cả các kết quả có thể xảy ra khi lấy 2 viên bi):}
			Số phần tử của không gian mẫu là $n(\Omega) = 5 \times 4 = 20$.
			Các phần tử của không gian mẫu là:
			\begin{itemize}
				\item Lấy Đ$_1$ trước: (Đ$_1$, Đ$_2$), (Đ$_1$, Đ$_3$), (Đ$_1$, X$_1$), (Đ$_1$, X$_2$)
				\item Lấy Đ$_2$ trước: (Đ$_2$, Đ$_1$), (Đ$_2$, Đ$_3$), (Đ$_2$, X$_1$), (Đ$_2$, X$_2$)
				\item Lấy Đ$_3$ trước: (Đ$_3$, Đ$_1$), (Đ$_3$, Đ$_2$), (Đ$_3$, X$_1$), (Đ$_3$, X$_2$)
				\item Lấy X$_1$ trước: (X$_1$, Đ$_1$), (X$_1$, Đ$_2$), (X$_1$, Đ$_3$), (X$_1$, X$_2$)
				\item Lấy X$_2$ trước: (X$_2$, Đ$_1$), (X$_2$, Đ$_2$), (X$_2$, Đ$_3$), (X$_2$, X$_1$)
			\end{itemize}
			Tổng cộng có $4+4+4+4+4 = 20$ kết quả.
			
			\item \textbf{Biến cố A: "lấy được 2 viên bi xanh" (các kết quả thuận lợi):}
			Các kết quả thuận lợi cho biến cố A là các cặp mà cả hai viên bi đều là bi xanh.
			Từ danh sách không gian mẫu ở trên, ta chọn ra các cặp thỏa mãn:
			\begin{itemize}
				\item (X$_1$, X$_2$)
				\item (X$_2$, X$_1$)
			\end{itemize}
			Vậy, số kết quả thuận lợi cho biến cố A là $n(A) = 2$.
			
			\item \textbf{Tính xác suất:}
			Xác suất để lấy được 2 viên bi xanh là:
			$$ P(A) = \frac{n(A)}{n(\Omega)} = \frac{2}{20} = \frac{1}{10} $$
		\end{itemize}
		\vspace{0.5cm}
		Kết luận: Xác suất để lấy được 2 viên bi là bi xanh là $\dfrac{1}{10}$.
	}
\end{ex}



\begin{ex}
	Một nhà may A theo kế hoạch phải may 200 bộ quần áo đồng phục cho học sinh khối 6 của một trường THCS theo thời gian quy định (mỗi ngày nhà may A may được số bộ quần áo là như nhau). Khi thực hiện, do cải tiến kỹ thuật nên mỗi ngày nhà may A may thêm được 5 bộ quần áo và hoàn thành kế hoạch trước 2 ngày. Hỏi theo kế hoạch, mỗi ngày nhà may A may được bao nhiêu bộ quần áo đồng phục?
	\loigiai{
		Gọi số bộ quần áo đồng phục nhà may A dự định may mỗi ngày theo kế hoạch là $x$ (bộ).
		Điều kiện: $x \in \mathbb{N}^*$ (x là số nguyên dương).
		
		Thời gian dự định để hoàn thành 200 bộ quần áo theo kế hoạch là: $\dfrac{200}{x}$ (ngày).
		
		Khi thực hiện, do cải tiến kỹ thuật, mỗi ngày nhà may A may được: $x+5$ (bộ).
		Thời gian thực tế để hoàn thành 200 bộ quần áo là: $\dfrac{200}{x+5}$ (ngày).
		
		Vì nhà may hoàn thành kế hoạch trước 2 ngày so với thời gian quy định, nên ta có phương trình:
		$$ \dfrac{200}{x} - \dfrac{200}{x+5} = 2 $$
		Để giải phương trình, ta quy đồng mẫu số:
		\begin{align*}
			\frac{200(x+5)}{x(x+5)} - \frac{200x}{x(x+5)} &= \frac{2x(x+5)}{x(x+5)} \\
			200(x+5) - 200x &= 2x(x+5) \\
			200x + 1000 - 200x &= 2x^2 + 10x \\
			1000 &= 2x^2 + 10x \\
			2x^2 + 10x - 1000 &= 0
		\end{align*}
		Chia cả hai vế cho 2, ta được phương trình bậc hai:
		$$ x^2 + 5x - 500 = 0 $$
		Tính biệt thức $\Delta$:
		$\Delta = b^2 - 4ac = 5^2 - 4(1)(-500) = 25 + 2000 = 2025$.
		Vì $\Delta > 0$, phương trình có hai nghiệm phân biệt:
		$\sqrt{\Delta} = \sqrt{2025} = 45$.
		$$ x_1 = \frac{-b + \sqrt{\Delta}}{2a} = \frac{-5 + 45}{2(1)} = \frac{40}{2} = 20 $$
		$$ x_2 = \frac{-b - \sqrt{\Delta}}{2a} = \frac{-5 - 45}{2(1)} = \frac{-50}{2} = -25 $$
		Đối chiếu với điều kiện $x \in \mathbb{N}^*$, ta nhận nghiệm $x=20$ và loại nghiệm $x=-25$.
		
		Vậy, theo kế hoạch, mỗi ngày nhà may A may được 20 bộ quần áo đồng phục.
	}
\end{ex}

\begin{ex}
	Cho $\triangle ABC$ có 3 góc nhọn. Vẽ đường tròn tâm O đường kính BC. Đường tròn (O) cắt AB; AC lần lượt ở D và E. BE và CD cắt nhau ở H.
	\begin{enumerate}
		\item[a.] Chứng minh bốn điểm A, D, H, E cùng nằm trên một đường tròn. Xác định vị trí tâm I của đường tròn đó.
		%\item[b.] Chứng minh: AE.AC=AB.AD 	Chứng minh: IE là tiếp tuyến của đường tròn (O)
	\end{enumerate}
	\loigiai{
		\begin{center}
			\begin{tikzpicture}[line join = round, line cap=round,>=stealth,font=\footnotesize,scale=1]
				% Def
				\def\r{3}
				
				% Coordinates
				\coordinate (O) at (0,0);
				\coordinate (B) at (-\r,0);
				\coordinate (C) at (\r,0);
				\coordinate (A) at (-2,5);
				
				\coordinate (D) at ($(A)!(C)!(B)$); % D la hinh chieu cua C len AB
				\coordinate (E) at ($(A)!(B)!(C)$); % E la hinh chieu cua B len AC
				
				% De tim H, can \usetikzlibrary{intersections} o preamble
				\path[name path=DC] (D) -- (C);
				\path[name path=BE] (B) -- (E);
				\coordinate[name intersections={of=DC and BE, by=H}]; % H la giao diem DC va BE
				
				\coordinate (I) at ($(A)!0.5!(H)$); % I la trung diem cua BC (gia dinh)
				
				% Draw
				\draw (O) circle (\r); % Ve duong tron tam O ban kinh r
				\draw (A) -- (B) -- (C) -- cycle; % Ve tam giac ABC
				\draw (C) -- (D); % Ve duong cao CD (tu C xuong AB)
				\draw (B) -- (E); % Ve duong cao BE (tu B xuong AC)
				\draw (A) -- (H);
				
				% Ky hieu goc vuông bằng lệnh %\pic[draw, angle radius=3mm]{right angle = C--A--B};
				% Vi du:
				\pic[draw, angle radius=2mm]{right angle = C--D--A}; % Goc CDA vuong
				\pic[draw, angle radius=2mm]{right angle = B--E--A}; % Goc BEA vuong
				
				% Fill 
				\fill (O) circle (1.5pt);
				\fill (B) circle (1.5pt);
				\fill (C) circle (1.5pt);
				\fill (A) circle (1.5pt);
				\fill (D) circle (1.5pt);
				\fill (E) circle (1.5pt);
				\fill (H) circle (1.5pt);
				\fill (I) circle (1.5pt);
				
				% Nodes	
				\node[below left] at (O) {$O$};
				\node[below left] at (B) {$B$};
				\node[below right] at (C) {$C$};
				\node[above] at (A) {$A$};
				\node[above left] at (D) {$D$}; % Vi tri nhan co the can dieu chinh
				\node[above right] at (E) {$E$}; % Vi tri nhan co the can dieu chinh
				\node[below] at (H) {$H$}; % Vi tri nhan co the can dieu chinh
				\node[left] at (I) {$I$}; % Vi tri nhan co the can dieu chinh
				% Ve va danh dau doan thẳng bằng nhau
			\end{tikzpicture}
			
		\end{center}
		\begin{enumerate}
			\item[a.] Chứng minh bốn điểm A, D, H, E cùng nằm trên một đường tròn. Xác định vị trí tâm I của đường tròn đó.
			\item[b.] Chứng minh: AE.AC=AB.AD 	Chứng minh: IE là tiếp tuyến của đường tròn (O)
			
			\textbf{Chứng minh $AE.AC=AB.AD$}
			
			Xét $\triangle AEB$ và $\triangle ADC$
			
			$\widehat{ADH} = \widehat{AEB} = 90^\circ$
			$\widehat{A}$ chung
			
			$\Rightarrow \triangle AEB \sim \triangle ADC$ (góc - góc)
			$\Rightarrow AE \cdot AC = AB \cdot AD$
			
		
			\textbf{IE là tiếp tuyến của đường tròn (O)}
			
			Ta có $IA=IH \Rightarrow IA=IE=IH = \frac{1}{2}AH$ (tính chất trung tuyến của tam giác vuông) $\Rightarrow \triangle IAE$ cân ở I $\Rightarrow \widehat{IEA} = \widehat{IAE}$.
			
			Mà $\widehat{IAE} = \widehat{EBC}$ (cùng phụ với góc ECB) và $\widehat{AEI} = \widehat{xEC}$ (đối đỉnh). Do $\triangle OEC$ cân ở O $\Rightarrow \widehat{OEC} = \widehat{OCE}$
			
			$\Rightarrow \widehat{xEC} + \widehat{CEO} = \widehat{EBC} + \widehat{ECB} = 90^\circ$. Hay $\widehat{xEO} = 90^\circ$
			
			Vậy $OE \perp IE$ tại điểm E nằm trên đường tròn (O) $\Rightarrow$ đpcm.
			
		\end{enumerate}
		}
\end{ex}
